
\section*{3. Construction of quotients \(F\backslash G\)
  (for \(G,F\) of finite type)}
\phantomsection\label{sec:sga3-VIA-3}
\addcontentsline{toc}{section}{3. Construction of quotients
  \texorpdfstring{\(F\backslash G\)}{F backslash G}
  (for G, F of finite type)}

\medskip
\phantomsection\label{par:sga3-VIA-3.1}%
\noindent\textbf{3.1.---}
Let \(A\) be an Artinian local ring, and let
\(u:F\to G\) be a homomorphism of \(A\)-groups. If
\[
  \mu:F\times_A F\longrightarrow F
  \qquad\text{and}\qquad
  \nu:G\times_A G\longrightarrow G
\]
denote the multiplication morphisms, and if \(\lambda\) denotes the
composite morphism
\begin{center}
  \phantomsection\label{diag:sga3-VIA-3.1-lambda-composite}%
  \begin{tikzcd}[column sep=large]
    F\times_A G
      \arrow[r,"u\times G"]
    &
    G\times_A G
      \arrow[r,"\nu"]
    &
    G
  \end{tikzcd},
\end{center}
recall that the \emph{left quotient} \(F\backslash G\) of \(G\) by
\(F\) is the cokernel of the \((\mathrm{Sch}/A)\)-groupoid \(G_*\)
described below:
\begin{center}
  \phantomsection\label{diag:sga3-VIA-3.1-groupoid}%
  \begin{tikzcd}[column sep=huge]
    F\times_A F\times_A G
      \arrow[r,shift left=2.1ex,"F\times\lambda"]
      \arrow[r,"\mu\times G" description]
      \arrow[r,shift right=2.1ex,"\operatorname{pr}_{2,3}"']
    &
    F\times_A G
      \arrow[r,shift left=1.2ex,"\lambda"]
      \arrow[r,shift right=1.2ex,"\operatorname{pr}_2"']
    &
    G
  \end{tikzcd}
\end{center}
Here \(\operatorname{pr}_2\) and \(\operatorname{pr}_{2,3}\) are the
projections of \(F\times_A G\) and
\(F\times_A(F\times_A G)\), respectively, onto the second factors.
We shall say that \(G_*\) is the \emph{groupoid with base \(G\)
defined by \(u\)} (cf.\ Expos\'e~V, \S~2.a; as in Expos\'e~V, we do
not follow in this expos\'e the convention of IV, \SGARefLink{sga3:iv:scholium:4:6:15}{4.6.15}).

Since the unique \(A\)-morphism \(F\to\operatorname{Spec}A\) is
universally open (EGA~IV\textsubscript{2}, 2.4.9),
\(\operatorname{pr}_2\) is an open morphism. The same therefore holds
for \(\lambda\), which is the composite of \(\operatorname{pr}_2\)
with the automorphism \(\sigma\) of \(F\times_A G\) defined by
\[
  \sigma(S)(x,y)=\bigl(x,u(S)(x)\cdot y\bigr),
\]
where \(S\) is a variable \(A\)-scheme and
\(x\in F(S)\), \(y\in G(S)\). One sees in the same way that
\(\operatorname{pr}_2\) and \(\lambda\) are flat when \(F\) is flat
over \(A\).

Let us also note, to conclude these preliminaries, that every
\(A\)-morphism \(s:\operatorname{Spec}A\to G\) defines an automorphism
of the groupoid \(G_*\) which induces on \(G\), \(F\times_A G\), and
\(F\times_A F\times_A G\), respectively, the automorphisms
\[
  r_s,\qquad
  \mathrm{id}_F\times_A r_s,\qquad
  \mathrm{id}_F\times_A\mathrm{id}_F\times_A r_s.
\]
We shall again write \(r_s\) for this automorphism of \(G_*\), and
shall say that \(r_s\) is the \emph{right translation defined by
\(s\)} (cf.\ \SGARefLink{sga3:VIA:0.4}{0.4}).

\medskip
\phantomsection\label{thm:sga3-VIA-3.2}%
\noindent\textbf{Theorem 3.2.}
\textit{Let \(F\) and \(G\) be flat groups locally of finite type over
an Artinian local ring \(A\). Let \(u:F\to G\) be a quasi-compact
homomorphism of \(A\)-groups whose kernel is finite over \(A\). Then:}%
\footnote{\label{note:sga3-VIA-33}Editor's note: we have added
(ii\(^{\prime}\)) and expanded point~\hyperref[concl:sga3-VIA-3.2-iv]{(iv)}, taking account of the
additions made in \SGARefLink{prop:sga3-VIA-2.5.2}{2.5.2}, \SGARefLink{prop:sga3-VIA-2.5.4}{2.5.4}, and \SGARefLink{lem:sga3-V-6.1}{Expos\'e~V, 6.1}.}
\begin{enumerate}
\item[\SGAInlineRefTarget{thm:sga3-VIA-3.2:item:i}(i)] \textit{The left quotient \(F\backslash G\) of \(G\) by
\(F\) exists in \((\mathrm{Sch}/A)\), and the sequence}
\begin{center}
  \phantomsection\label{diag:sga3-VIA-3.2-exact-sequence}%
  \begin{tikzcd}[column sep=large]
    F\times_A G
      \arrow[r,shift left=1.2ex,"\lambda"]
      \arrow[r,shift right=1.2ex,"\operatorname{pr}_2"']
    &
    G
      \arrow[r,"p"]
    &
    F\backslash G
  \end{tikzcd}
\end{center}
\textit{is exact in the category of all ringed spaces.}

\item[\SGAInlineRefTarget{thm:sga3-VIA-3.2:item:ii}(ii)] \textit{The canonical morphism
\(p:G\to F\backslash G\) is surjective and open, and
\(F\backslash G\) is a direct sum of schemes of finite type over
\(A\).}

\item[(ii\(^{\prime}\))]\phantomsection\label{concl:sga3-VIA-3.2-ii-prime} \textit{More precisely,
\(X=F\backslash G\) is equipped with a right action of \(G\) such
that \(p(e)\cdot g=p(g)\) for every \(g\in G\). Consequently, the
connected components of \(X\) are of finite type over \(A\),
irreducible, and all of dimension \(\dim G-\dim F\).}

\item[\SGAInlineRefTarget{thm:sga3-VIA-3.2:item:iii}(iii)] \textit{The canonical morphism
\[
  F\times_A G\longrightarrow G\times_{(F\backslash G)}G
\]
is surjective.}

\item[(iv)]\phantomsection\label{concl:sga3-VIA-3.2-iv} \textit{If \(u\) is a monomorphism,}%
\footnote{\label{note:sga3-VIA-34}Editor's note: in this case, the
hypothesis that \(G\) be flat can be omitted; cf.\ \SGARefLink{par:sga3-VIA-3.3}{subsection~3.3}.}
\textit{ then:}
\begin{enumerate}
\item[\SGAInlineRefTarget{thm:sga3-VIA-3.2:item:a}(a)] \textit{The morphism
\[
  F\times_A G
    \xrightarrow{(\lambda,\operatorname{pr}_2)}
  G\times_{(F\backslash G)}G
\]
is an isomorphism, and \(G\to F\backslash G\) is faithfully flat and
locally of finite presentation.}

\item[(a\(^{\prime}\))] \textit{\(F\backslash G\) represents the
\((\mathrm{fppf})\) quotient sheaf
\(\widetilde{F\backslash G}\), and \(G\to F\backslash G\) is an
\(F\)-torsor locally trivial for the \((\mathrm{fppf})\) topology.}

\item[(b)] \textit{\(F\backslash G\) is flat over \(A\), and is
smooth over \(A\) if \(G\) is so.}

\item[(c)] \textit{\(u:F\to G\) is a closed immersion, and
\(F\backslash G\) is separated.}

\item[(d)] \textit{If, moreover, \(F\) is a normal subgroup of
\(G\), there is a unique \(A\)-group structure on \(F\backslash G\)
for which \(p:G\to F\backslash G\) is a morphism of \(A\)-groups.}
\end{enumerate}
\end{enumerate}

In the proof of this theorem, \(A'\) will denote an \(A\)-algebra
which is local, finite, and free over \(A\). If \(R\) is a relation
involving \(A'\), we shall say that ``\(R(A')\) is true when \(A'\)
is sufficiently large'' if there is an \(A\)-algebra \(A_1\), local,
finite, and free over \(A\), such that the relation \(R(A')\) holds
for every \(A\)-algebra \(A'\) which is local, finite, and free over
\(A_1\).

We shall first prove the theorem when \(F\) and \(G\) are of finite
type over \(A\).
